\section{The \texttt{Sapphire\_Filter} McStas Component}
Sapphire filter at 300K

\subsection*{Identification}
\begin{itemize}
  \item \textbf{Author:} Jonas Okkels Birk (based upon Filteg\_Graphite by Thomas C Hansen (2000))
  \item \textbf{Origin:} PSI
  \item \textbf{Date:} 6 December 2006
\end{itemize}

\subsection*{Description}
ESCRIPTION

This sapphire filter, defined by two identical rectangular opening apertures, is based upon an absorption function absorp=exp(L*A+C*(exp(-(B/L\textasciicircum{}2+D/L\textasciicircum{}4)))) decribed by Freund (1983) Nucl. Instrum. Methods, A278, 379-401. The defaultvalues is for Sapphire at 300K as found by measurement by (Mildner \& Lamaze (1998) J. Appl. Cryst. 31,835-840 ( http://journals.iucr.org/j/issues/1998/06/00/hw0070/hw0070bdy.html\#BB4)). It is possble to ajust the formular to other materials or temperatures by typing in other parameters. The transmission in sapphire is only lightly demependt on temperature, but The formular is only cheked at wavelenths between 0.5 and 14 \AA{} (0.4 meV 0.3 eV). The filter is for example used in the Eiger beamline at PSI to cut of high energies.

\subsection*{Input parameters}
Parameters in \textbf{boldface} are required; the others are optional.

\begin{longtable}{p{0.22\textwidth}p{0.12\textwidth}p{0.46\textwidth}p{0.14\textwidth}}
\toprule
\textbf{Name} & \textbf{Unit} & \textbf{Description} & \textbf{Default} \\
\midrule
\endhead
xmin & m & Lower x bound & -0.16 \\
xmax & m & Upper x bound & 0.16 \\
ymin & m & Lower y bound & -0.16 \\
ymax & m & Upper y bound & 0.16 \\
len & m & Thickness of filter & 0.1 \\
A & m$^{-1}$ \AA{}$^{-1}$ &  & 0.8116 \\
B & \AA{}$^{2}$ &  & 0.1618 \\
C & m$^{-1}$ &  & 21.24 \\
D & \AA{}$^{4}$ &  & 0.1291 \\
\bottomrule
\end{longtable}

\subsection*{Links}
\begin{itemize}
  \item Component source code found in file \texttt{Sapphire\_Filter.comp}.
\end{itemize}
\IfFileExists{contrib/Sapphire_Filter_static.tex}{\input{contrib/Sapphire_Filter_static.tex}}{}